% ============================================================
% Discussion / Limitations — three-pillar framing.
% ============================================================
\section{Discussion}
\label{sec:discuss}

\paragraph{Why state-evolution, not action.}
The central design choice in \memmark{} is to watermark the memory
backend's internal state evolution rather than the visible action
trajectory. This matters because action-level attribution is only useful
while the original trajectory remains available. Once memory has left
the system, been compacted, or been migrated into another storage
layer, the trajectory is usually gone. In contrast, the state-evolution
path survives as part of the memory object itself, which is why the R3
setting in \S\ref{sec:rq3} is the paper's most stringent test: it asks
whether attribution still works when the only trusted artifact is the
snapshot.

\paragraph{What distribution preservation buys.}
Lemma~\ref{lem:dist-preserve} is a marginal guarantee, not a claim that
every individual decision is unchanged. Its significance is that the
watermark is inserted without distorting the backend's preference
distribution in expectation. Empirically, that is exactly what RQ1
shows: utility metrics remain within statistical noise of the
unwatermarked baseline, and the sign of the delta is not consistently
negative. The \texttt{signed-metadata-only} baseline is especially
useful here because it isolates the overhead of keyed selection from the
audit-trace machinery around it.

\paragraph{Why the verification hierarchy matters.}
R1, R2, and R3 are not just increasingly weak setups; they correspond
to distinct forensic realities. R1 is the easy case with a full audit
log. R2 captures retention loss, truncation, or partial deletion of the
external log. R3 is the snapshot-only regime that motivates the paper.
The smooth degradation in R2 and the full recovery in R3 together show
that the audit design degrades gracefully rather than failing at the
first sign of missing evidence.

\paragraph{Limitations and future work.}
The current implementation is validated on \amem{} and \graphiti{}, so
broadening backend coverage to systems such as MemOS or MemMachine is an
engineering extension rather than a conceptual one. We also assume that
the secret key is not exfiltrated; if an adversary can both rewrite the
snapshot and recover the key, any keyed watermark is forgeable.
Finally, while the current candidate enumeration works well in our
benchmarks, larger \(K\) may increase capacity at the cost of higher
prompting overhead and potentially lower output quality. Promising
follow-ups include online attestation for streaming agents and
tighter cross-session anchor chaining for long-running deployments.
